\chapter{水体卫生  【重点章节】}
\setcounter{chapter}{5}
\chapterintro{
本章为\keyword{重点考察章节}（6学时）。掌握水质的评价指标（物理、化学、微生物学性状指标）；掌握水体污染来源及各种水体的污染特点；掌握水体自净机制及氧垂曲线；掌握水体富营养化的概念及危害；掌握水体污染的危害；了解水环境标准和水体卫生防护措施。
}

\section{水资源的种类及卫生学特征}

\begin{defbox}
\textbf{水资源（Water Resources）}：全球水量中对人类生存、发展可用的水量，主要指逐年可得到更新的那部分淡水量。
\end{defbox}

\begin{table}[htbp]
\centering
\caption{天然水资源分类及卫生学特征}
\begin{tabularx}{\textwidth}{lXp{3.5cm}}
\hline
\textbf{类型} & \textbf{卫生学特征} & \textbf{主要问题} \\
\hline
\imp{降水} & 水质较好，矿物质含量低；地区分布极不平衡 & 水量无保证；受大气质量影响（如酸雨）；沿海地区降水含碘较多 \\
\imp{地表水} & 溶解土壤矿物质；水量充足；可分为封闭型（湖水、水库水等）和开放型（江河水） & 人为污染是影响水质最主要因素；有丰水期和枯水期变化 \\
\imp{地下水} & 水质清洁，矿物质含量较高；水量较稳定 & 一旦污染，自净极其困难；硬度较高 \\
\hline
\end{tabularx}
\end{table}

\begin{defbox}
\textbf{透水层}：由颗粒较大的砂、砾石组成，能渗水与存水。

\textbf{不透水层}：由颗粒细小致密的黏土层和岩石层构成。

\textbf{浅层地下水}：位于第一个不透水层之上。水质物理性状较好，细菌数较地表水少，但硬度增高，溶解氧减少。我国广大农村使用。

\textbf{深层地下水}：位于第一个不透水层之下。水质透明无色，水温恒定，水量稳定，细菌数很少，但盐类含量高、硬度大。城镇或企业集中式供水常用。

\textbf{泉水}：地下水通过地表缝隙自行涌出。分潜水泉（浅层地下水涌出）和自流泉（深层地下水由裂隙中涌出）。
\end{defbox}

\section{水质的性状和评价指标}

\subsection{物理性状指标}

\begin{table}[htbp]
\centering
\caption{水质物理性状指标}
\begin{tabularx}{\textwidth}{lX}
\hline
\textbf{指标} & \textbf{含义及意义} \\
\hline
水温 & 影响水中生物和化学反应速率 \\
色 & 纯净水无色；有色可能提示污染 \\
臭和味 & 异常臭和味提示有机物分解或化学污染 \\
\imp{浑浊度} & 表示水中悬浮物和胶体物对光线透过时的阻碍程度。1度=1L水中含1mg标准硅藻土 \\
放射性 & 天然和人为放射性物质含量 \\
\hline
\end{tabularx}
\end{table}

\subsection{化学性状指标}

\begin{keybox}
\textbf{主要化学性状指标：}

\begin{enumerate}[leftmargin=*]
    \item \imp{pH值}：天然水pH一般在7.2--8.5之间
    \item \imp{总固体（Total Solid）}：水样蒸发至干后的残留物总量 = 溶解性固体 + 悬浮性固体
    \item \imp{硬度（Hardness）}：溶于水中钙、镁盐类的总含量，以CaCO$_3$（mg/L）表示
    \begin{itemize}
        \item 暂时硬度：水煮沸后能去除的硬度（碳酸盐硬度）
        \item 永久硬度：煮沸后不能去除的硬度（非碳酸盐硬度）
    \end{itemize}
    \item \imp{含氮化合物}——反映有机污染及其自净进程：
    \begin{itemize}
        \item \textbf{有机氮和蛋白氮}：有机氮是有机含氮化合物的总称，蛋白氮是已分解成较简单的有机氮。两者显著增高说明\keyword{水体新近受明显有机性污染}
        \item \textbf{氨氮}：含氮有机物在有氧条件下经微生物分解形成的最初产物，表示\keyword{新近可能有人畜粪便污染}
        \item \textbf{亚硝酸盐氮}：氨在有氧条件下经亚硝酸菌作用形成，是硝化过程的中间产物，含量高说明\keyword{有机物无机化过程尚未完成，污染仍存在}
        \item \textbf{硝酸盐氮}：含氮有机物氧化分解的最终产物，含量高而氨氮、亚硝酸盐氮不高，说明\keyword{过去曾受有机污染，现已完成自净}
    \end{itemize}
    \begin{table}[H]
    \centering
    \caption{三氮组合判断水体污染状况}
    \begin{tabular}{ccc}
    \hline
    \textbf{NH$_3$-N} & \textbf{NO$_2^-$-N / NO$_3^-$-N} & \textbf{污染状况} \\
    \hline
    高 & 低 & \textcolor{keyred}{新近污染} \\
    低 & 高 & \textcolor{darkgreen}{陈旧污染（已自净）} \\
    高 & 高 & 持续污染（过去和新近均有污染） \\
    均高 & — & 过去受污染，目前自净正在进行 \\
    \hline
    \end{tabular}
    \end{table}
    \item \imp{溶解氧（DO）}：溶解在水中的氧含量，清洁水DO较高
    \item \imp{化学耗氧量（COD）}：用强氧化剂氧化水中有机物所消耗的氧量
    \item \imp{生化需氧量（BOD）}：水中有机物在需氧微生物分解时消耗的溶解氧量。\textbf{BOD$_5^{20}$}：20°C培养5日后1L水中减少的溶解氧量。清洁水BOD$_5$一般 $<$ 1 mg/L
    \item \imp{总有机碳（TOC）}：水中全部有机物的含碳量，是评价水体有机物污染程度的综合性指标之一，但不能说明有机污染物的性质
    \item \imp{总需氧量（TOD）}：1L水中还原物质（有机物和无机物）在一定条件下被氧化时所消耗氧的毫升数，数值愈大污染愈严重。TOC和TOD的检测可实现快速自动化
    \item \imp{氯化物}：含量突然增高时表明水可能受到人畜粪便、生活污水或工业废水的污染
    \item \imp{硫酸盐}：含量突然增加表明水可能受到生活污水、工业废水或硫酸铵化肥等污染
\end{enumerate}
\end{keybox}

\subsection{微生物学性状指标}

\begin{defbox}
\textbf{指示菌（Indicator Bacteria）}：具有代表微生物污染总体状况的菌种。
\end{defbox}

\begin{enumerate}[leftmargin=*]
    \item \imp{细菌总数（Bacteria Count）}：1 mL水在普通琼脂培养基中37°C培养24小时后生长的细菌菌落数。
    \item \imp{总大肠菌群（Total Coliforms）}：需氧及兼性厌氧、37°C生长时使乳糖发酵、24小时内产酸产气的革兰氏阴性无芽孢杆菌。是\keyword{粪便污染指示菌}。
    \item \imp{粪大肠菌群（Fecal Coliforms）}：44.5°C培养仍能生长的大肠菌群，更特异地指示粪便污染。
\end{enumerate}

\section{水体的污染来源及特征}

\subsection{污染来源}

\begin{enumerate}[leftmargin=*]
    \item \imp{工业废水}：排放量大、成分复杂、毒性强、难处理
    \item \imp{生活污水}：含有机物、病原微生物、氮磷等
    \item \imp{农业污水}：含农药、化肥、畜禽养殖废弃物
\end{enumerate}

\subsection{各种水体的污染特点}

\begin{keybox}
\begin{enumerate}[leftmargin=*]
    \item \imp{河流}：污染程度取决于\keyword{径污比}（径流量与排入污水量的比值）。
    \item \imp{湖泊、水库}：稀释混合能力较差、水交换缓慢。主要问题是\keyword{富营养化}。
    \item \imp{地下水}：污染过程缓慢，不易发现；自净能力极差；微生物含量少；一旦污染恢复极其困难。
    \item \imp{海洋}：污染源多而复杂；污染范围大；持续性长。
\end{enumerate}
\end{keybox}

\section{水体污染自净和污染物的转归}

\subsection{水体自净}

\begin{defbox}
\textbf{水体自净（Self-Purification）}：水体受污染后，污染物在水体的物理、化学和生物学作用下，不断稀释、扩散、分解破坏或沉入水底，水中污染物浓度逐渐降低，水质最终恢复到污染前状况的过程。
\end{defbox}

\textbf{水体自净三阶段：}
\begin{enumerate}[leftmargin=*]
    \item 第一阶段：易被氧化的有机物进行化学氧化分解——数小时完成
    \item 第二阶段：有机物在微生物作用下进行生物化学氧化分解——一般延续数日
    \item 第三阶段：含氮有机物的硝化过程——一般延续一个月左右
\end{enumerate}

\textbf{水体自净过程的特征：}污染物浓度逐渐下降；溶解氧（DO）先降后升；生物种群随之变化；水质逐渐恢复。

\textbf{环境容纳量（Environmental Capacity）}：天然水体能接纳一定量的污染物进行自净，使水质成分保持平衡的能力。

\subsection{水体自净的机制}

\begin{keybox}
\textbf{三种净化机制：}

\begin{enumerate}[leftmargin=*]
    \item \imp{物理净化}：稀释、混合、吸附、沉淀等
    \item \imp{化学净化}：氧化还原、水解、光解等
    \item \imp{生物净化}（最为重要最为活跃）：生物通过代谢作用分解水中污染物。包括：
    \begin{itemize}
        \item 需氧微生物将有机物分解为CO$_2$、水、硫酸盐、硝酸盐等
        \item 病原体的死灭
        \item 藻类对氮磷的吸收
        \item \imp{水生植物净化}：浮萍、凤眼莲、芦苇等能吸收、分解或浓缩水中汞、镉、锌等重金属及难降解有机物。例如水葫芦可吸收N、P，芦苇能分解酚类，浮萍能富集镉
    \end{itemize}
\end{enumerate}
\end{keybox}

\subsection{氧垂曲线}

\begin{keybox}
\textbf{氧垂曲线（Oxygen Sag Curve）}反映有机污染物排入河流后溶解氧的变化：

\begin{center}
\resizebox{0.95\textwidth}{!}{%
\begin{tikzpicture}[yscale=0.85, xscale=0.95]
    % Axes
    \draw[->, thick] (0,0) -- (12,0) node[below] {距离（沿河流方向）};
    \draw[->, thick] (0,0) -- (0,6.5) node[above] {溶解氧 DO (mg/L)};

    % Saturated DO line
    \draw[dashed, gray!60, thin] (0,5) -- (12,5);
    \node[gray!60, anchor=south west] at (12,5) {\footnotesize\ 饱和溶解氧};

    % DO curve
    \draw[thick, main] (0,5)
        .. controls (2,4.6) and (3,3.2) .. (4.5,2.5)
        .. controls (6,1.8) and (7.5,3.2) .. (10,4.9)
        .. controls (11,5.5) and (11.8,5.3) .. (12,5.2);
    \node[main, anchor=west] at (12,5.2) {\footnotesize\ DO曲线};

    % Critical point
    \draw[fill=keyred] (4.5,2.5) circle (2.5pt);
    \node[keyred, anchor=north west] at (4.3,2.2) {\small $C_p$（最大缺氧点）};

    % Health standard line
    \draw[red!60, dashed, thick] (0,1.5) -- (12,1.5);
    \node[red!60, anchor=south west] at (12,1.5) {\footnotesize\ 卫生标准 (4 mg/L)};

    % Zone labels
    \draw[gray!40, thick, decorate, decoration={snake, amplitude=0.5mm}] (0,5.8) -- (2.5,5.8);
    \node[gray!60] at (1.3,6.1) {\small 清洁区};

    \draw[gray!40, thick, decorate, decoration={snake, amplitude=0.5mm}] (2.5,5.8) -- (6.5,5.8);
    \node[gray!60] at (4.5,6.1) {\small 分解区（恶化区）};

    \draw[gray!40, thick, decorate, decoration={snake, amplitude=0.5mm}] (6.5,5.8) -- (12,5.8);
    \node[gray!60] at (9,6.1) {\small 恢复区};

    % Vertical dashed lines for zones
    \draw[gray!30, dashed] (2.5,0.5) -- (2.5,5.8);
    \draw[gray!30, dashed] (6.5,0.5) -- (6.5,5.8);
\end{tikzpicture}%
}
\end{center}

\begin{itemize}[leftmargin=*]
    \item \textbf{Cp点}：之前耗氧 $>$ 复氧，DO下降，水质恶化；之后复氧 $>$ 耗氧，DO逐渐恢复
    \item 若Cp点DO $>$ 4 mg/L：排放未超过水体自净能力
    \item 若Cp点DO $<$ 4 mg/L：水质严重恶化，可能出现黑臭
\end{itemize}
\end{keybox}

\subsection{水体污染物的转归}

\begin{itemize}[leftmargin=*]
    \item \imp{迁移}：污染物在水体中的空间位移
    \item \imp{生物富集作用}：生物从环境中摄入低浓度物质在体内逐渐累积
    \item \imp{生物放大作用}：污染物沿食物链逐级浓缩放大
\end{itemize}

\begin{tipbox}
\textbf{DDT在水生食物链中的迁移和转归}：水（0.00005 ppm）→ 浮游生物（0.04 ppm）→ 小鱼（0.5 ppm）→ 大鱼（2.0 ppm）→ 水鸟（25 ppm）。总体浓缩系数达50万倍。
\end{tipbox}

\section{水体污染的危害}

\subsection{生物性污染的危害}

\begin{enumerate}[leftmargin=*]
    \item \imp{介水传染病}：霍乱、伤寒、痢疾、甲型肝炎等
    \item \imp{藻类毒素污染}：
    \begin{itemize}
        \item 蓝藻水华产生\keyword{微囊藻毒素（MC）}，具有肝毒性和促癌作用。分为MC-LR（亮氨酸，急性毒性最强，抑制丝氨酸/苏氨酸蛋白磷酸酶-1和-2A活性）、MC-RR（精氨酸）、MC-YR（酪氨酸）。肝脏是MC的主要靶器官，可引起肝脏大面积肿胀出血、坏死、肝细胞结构功能破坏甚至肝衰竭
        \item MC耐热性强，煮沸不能破坏其毒性；常规水处理不能完全消除
        \item MC被确定为\imp{2B类致癌物}，是迄今已发现的最强\imp{肝癌促进剂}，且MC-LR和黄曲霉毒素B1有\imp{协同促癌作用}
        \item 慢性暴露可引起血清中谷丙转氨酶（ALT）、碱性磷酸酶（ALP）等多种酶水平显著改变
        \item 我国生活饮用水水质标准（2022）中MC-LR为扩展指标，限值1 $\mu$g/L
        \item \textbf{贝毒}：麻痹性贝毒（70\%中毒/死亡事件）、腹泻性贝毒（损伤肝脏）、记忆缺失性贝毒、神经性贝毒、西加鱼毒
    \end{itemize}
\end{enumerate}

\subsection{化学性污染的危害}

\begin{keybox}
\textbf{主要化学污染物及其危害：}

\begin{enumerate}[leftmargin=*]
    \item \imp{酚类化合物}：使水产生异臭（氯酚臭阈低至5 $\mu$g/L）；低浓度使蛋白质变性，高浓度使蛋白质沉淀；对皮肤黏膜有强烈刺激腐蚀作用；急性酚中毒表现为大量出汗、肺水肿、吞咽困难、肝及造血系统损害、黑尿
    \item \imp{多氯联苯（PCBs）}：无色或淡黄色油状有机氯化合物，性质稳定，耐酸碱、耐腐蚀。作为首批受控POPs，具持久性、蓄积性、迁移性、高毒性。主要经食物摄入暴露。\textbf{典型案例：日本米糠油中毒事件、中国台湾油症事件。}健康危害：①肝肿大和脂肪肝；②皮疹、痤疮、色素沉着，严重者肝损害、黄疸、肝性脑病；③雌激素样作用——拮抗雄激素睾酮、干扰雌激素正常代谢、干扰甲状腺功能；④通过乳汁和胎盘影响子代（胎儿油症：新生儿体重减轻、皮肤颜色异常）
    \item \imp{汞（Hg）}：转化为甲基汞后毒性大增；导致\keyword{水俣病}——慢性甲基汞中毒，以神经系统损害为主。甲基汞易于被水生生物吸收，通过生物富集和生物放大作用，使鱼贝类甲基汞富集百万倍以上
    \item \imp{邻苯二甲酸酯类（PAEs）}：塑料增塑剂、软化剂，易从塑料溶出。急性毒性很低但具有较强的雄性生殖毒性，是典型的EDCs，睾丸是重要靶器官。人类主要通过食物摄入暴露。DBP限值0.003 mg/L，DEHP限值0.008 mg/L（DEHP为2B类致癌物）
    \item \imp{全氟化合物（PFASs）}：持久性有机污染物
\end{enumerate}
\end{keybox}

\subsection{物理性污染的危害}

\begin{itemize}[leftmargin=*]
    \item \imp{热污染}：工业冷却水导致水温升高→DO降低→影响水生生态系统→促进藻类繁殖加剧富营养化。热污染还可能引起致病微生物的孳生繁殖，如变形原虫在发电厂废热水中孳生繁衍，进入人体后引起脑膜炎
    \item \imp{放射性污染}：核电站事故、核废料排放等。人体接触高浓度放射性物质的水可引起外照射，饮水或食品受放射性污染后可造成内照射
\end{itemize}

\section{水环境标准}

\subsection{地表水环境质量标准}

\textbf{制订地表水环境质量标准的原则：}
\begin{enumerate}[leftmargin=*]
    \item 防止地表水被污染而传播疾病
    \item 防止地表水被污染而危害人体健康
    \item 保证地表水感官性状良好
    \item 保证地表水自净过程正常进行
\end{enumerate}

\subsection{废水处理分级}

\begin{table}[htbp]
\centering
\caption{废水处理三级体系}
\begin{tabularx}{\textwidth}{lX}
\hline
\textbf{级别} & \textbf{处理方法与去除对象} \\
\hline
一级处理 & 物理方法——去除悬浮固体、漂浮物 \\
二级处理 & \imp{生物处理方法}——去除溶解性和胶体有机物（BOD去除） \\
三级处理 & 深度处理——去除氮磷、重金属、难降解有机物等 \\
\hline
\end{tabularx}
\end{table}

\begin{defbox}
\textbf{中水回用}：将生活污水和工业废水经处理后达到一定水质标准，回用于非饮用用途（如绿化、冲厕、工业冷却等）。
\end{defbox}

\section{水体污染的预防对策}

\begin{itemize}[leftmargin=*]
    \item \imp{污染源头控制}——污染物尚未对水体造成污染之前采用积极有效的措施，防止污染物进入水体
    \item \imp{清洁生产（Cleaner Production）}——将综合预防的环境保护策略持续应用于生产过程和产品中，以减少污染物对人类和环境的风险。内容：清洁的能源、清洁的生产过程、清洁的产品
    \item 工业废水的利用与处理
    \item 生活污水的利用与处理
    \item 医院污水的处理（含病原体、放射性物质等）
    \item 水体污染的卫生调查、监测和监督
\end{itemize}

\section{水体污染监测要点}

\subsection{江河水系监测}
\begin{enumerate}[leftmargin=*]
    \item \textbf{采样断面：}至少3个——①清洁/对照断面（污染源上游）；②污染断面（污染源下游）；③自净断面（污染断面下游至少1.5 km以上，水体基本达到自净的地方）。还需考虑表层、中层和底层的垂直布点
    \item \textbf{采样点：}水量不大/河面较窄→单点（河中心）；水量大/河面宽→多点（3--5个点，离岸50 m、150 m、江心）。采表层水应在水面下20--50 cm处
    \item \textbf{采样时间：}丰水期、枯水期和平水期各采样一次，每次连续2--3天，避开雨天。有潮汐河流应分别在高潮和低潮时采样
    \item \textbf{必测项目：}水温、浑浊度、色度、pH、总硬度、DO、BOD、总大肠菌群、挥发性酚、氰化物、砷、汞、铬。五项毒物必测：酚、氰、汞、砷、六价铬
\end{enumerate}

\subsection{底质与生物监测}
\begin{itemize}[leftmargin=*]
    \item \imp{底质监测}：底质中有害物质含量的垂直分布能反映水体污染的历史状况；某些污染物在水中含量很低不易检出，而在底质中含量可比水中高出很多倍。项目：汞、铬、砷、氰化物、酚类等五项毒物及Cu、Cd、Zn、Mn等
    \item \imp{水生生物监测}：生物种群/数量/分布测定；生物体内毒物负荷测定；污染物对水生生物综合作用检测；大肠菌群和病原微生物检测
\end{itemize}

\subsection{其他水体监测}
\begin{itemize}[leftmargin=*]
    \item \imp{湖泊水库}：增加磷（P）、氮（N）及藻类毒素的测定
    \item \imp{地下水}：增测碘、氟、砷、硫化物和硝酸盐等
\end{itemize}

\subsection{水体污染调查类型}
\begin{enumerate}[leftmargin=*]
    \item 基础调查——了解水体基本状况，范围较大
    \item 监测性调查——定期对水污染进行监测
    \item 专题调查——针对特定污染类型或污染物的专门调查
    \item 应急性突发事件调查——在严重污染事故时进行
\end{enumerate}

\section{复习题}

\begin{enumerate}[leftmargin=*, itemsep=8pt]
    \item \textbf{【名词解释】}水体自净；富营养化；水华；赤潮
    \item \textbf{【名词解释】}DO；COD；BOD$_5$；TOC；TOD
    \item \textbf{【名词解释】}硬度（暂时硬度与永久硬度）；清洁生产
    \item \textbf{【名词解释】}总大肠菌群；细菌总数
    \item \textbf{【名词解释】}氧垂曲线；生物放大作用
    \item \textbf{【简答题】}简述水体自净过程的特征及机制（含物理、化学、生物净化）。
    \item \textbf{【简答题】}简述地下水污染的特点。
    \item \textbf{【简答题】}氧垂曲线中Cp点的含义是什么？其卫生学意义何在？
    \item \textbf{【简答题】}简述水体富营养化的形成过程及危害。
    \item \textbf{【简答题】}请根据含氮化合物指标判断水体的污染状况。
    \item \textbf{【简答题】}简述制订地表水环境质量标准的原则。
    \item \textbf{【简答题】}废水处理分为哪三级？
    \item \textbf{【简答题】}简述多氯联苯（PCBs）的环境及健康危害。
    \item \textbf{【非标准答案题】}某水产养殖水体在施撒有机肥后湖面不断冒出气泡，短期内突然有大量鱼类死亡，请分析最可能的原因及应对措施。
\end{enumerate}

\subsection*{参考答案要点}

\begin{enumerate}[leftmargin=*, itemsep=5pt]
    \item \textbf{水体自净}：水体受污染后，污染物在物理、化学和生物学作用下浓度逐渐降低、水质恢复的过程。\textbf{富营养化}：湖泊水库接纳过多含P、N污水时藻类大量繁殖的现象。淡水中称\imp{水华}，海湾中称\imp{赤潮}。

    \item \textbf{DO}：溶解氧。\textbf{COD}：化学耗氧量，用强氧化剂氧化有机物消耗的氧量。\textbf{BOD$_5$}：20°C培养5日，微生物分解有机物消耗的溶解氧量。\textbf{TOC}：总有机碳。

    \item \textbf{硬度}：水中钙、镁盐类总含量（CaCO$_3$ mg/L）。暂时硬度：煮沸可去除（碳酸盐）。永久硬度：煮沸不能去除（非碳酸盐）。

    \item \textbf{总大肠菌群}：需氧及兼性厌氧、37°C使乳糖发酵产酸产气的革兰氏阴性无芽孢杆菌，粪便污染指示菌。\textbf{细菌总数}：1mL水37°C培养24h的菌落数。

    \item \textbf{氧垂曲线}：反映有机污染物排入河流后DO沿程变化的曲线。\textbf{生物放大作用}：重金属/难降解有机物通过食物链逐级浓缩放大。

    \item 特征：污染物浓度下降、DO先降后升、生物种群变化、水质恢复。机制：物理净化（稀释混合吸附沉淀）、化学净化（氧化还原水解）、生物净化（有机物降解、病原体死灭）。

    \item 污染过程缓慢不易发现；流动极缓慢自净能力差；微生物含量少；一旦污染恢复极其困难。

    \item Cp点为溶解氧最低点。之前耗氧$>$复氧，之后复氧$>$耗氧。若DO$>$4mg/L则未超自净能力；若$<$4mg/L则水质恶化。

    \item 形成：P、N污水排入→藻类大量繁殖→DO降低→水质恶化→水生生物死亡。危害：水质恶化、藻毒素产生、影响供水、破坏生态系统。

    \item NH$_3$高+NO$_3^-$低=新近污染；NH$_3$低+NO$_3^-$高=陈旧已自净；均高=持续污染。

    \item 四项原则：防止传播疾病；防止危害健康；保证感官性状良好；保证自净过程正常进行。

    \item 一级（物理去悬浮物）、二级（生物去有机物）、三级（深度处理去N、P、重金属等）。

    \item PCBs持久性、蓄积性强，具免疫、生殖、神经毒性和致癌性；可远距离迁移；米糠油事件元凶。

    \item 有机肥分解消耗大量DO→水体缺氧→鱼类窒息死亡。对策：控制施肥量、增加曝气、监测DO。
\end{enumerate}

% ----- PPT参考题 -----
\subsection*{参考题（PPT课件补充）}

\begin{enumerate}[leftmargin=*, itemsep=8pt]
    \item \textbf{【单选题】}下列水体中，卫生学特征符合"水量大、水质好"的水资源是（~~~~）
    \begin{multicols}{5}
    A. 雨雪水 \\
    B. 湖泊 \\
    C. 地下水 \\
    D. 河流 \\
    E. 海水
    \end{multicols}
    \item \textbf{【单选题】}当水中未检出亚硝酸盐氮，但氨氮和硝酸盐氮含量均明显增高时，可推测该水体受有机污染的状况是（~~~~）
    \begin{multicols}{3}
    A. 刚刚发生污染 \\
    B. 受污染不久，现仍在继续 \\
    C. 长久以来一直受污染 \\
    D. 旧污染已分解，又有新的污染 \\
    E. 已完全自净
    \end{multicols}
    \item \textbf{【单选题】}水体富营养化更容易发生在以下哪个地方？（~~~~）
    \begin{multicols}{5}
    A. 嘉陵江 \\
    B. 太平洋 \\
    C. 滇池 \\
    D. 长江 \\
    E. 亚马逊河
    \end{multicols}
\end{enumerate}

\subsubsection*{参考答案}
\begin{enumerate}[leftmargin=*, itemsep=5pt]
    \item \textbf{C（地下水）}。地下水水质清洁，矿物质含量较高，水量较稳定，符合"水量大、水质好"的特征。雨雪水水量无保证；湖泊和河流属于地表水，易受人为污染；海水为咸水，非直接可用水资源。
    \item \textbf{D（旧污染已分解，又有新的污染）}。氨氮高提示有机物正在进行氨化（新污染），硝酸盐氮高提示硝化作用已完成（旧污染已自净）。两者同时增高说明旧污染已分解且又发生了新污染。
    \item \textbf{C（滇池）}。富营养化主要发生在流动性差、水交换缓慢的封闭/半封闭水体（如湖泊、水库）。滇池属于此类湖泊，而嘉陵江、长江、亚马逊河为流动性好的河流，太平洋为开放性海域。
\end{enumerate}

\newpage

% =====================================================================
%  CHAPTER 6: 饮用水卫生
% =====================================================================
\newpage